
\section{Analysis}

\subsection{What Does the Score Network Learn?}

\textbf{Figure 8.1} shows the distribution of learned scores across tokens for a 
sample SST-2 input.

\subsubsection{Content vs Function Words}

\textbf{Table 8.1: Mean score by POS tag.}

\begin{table}[H]
\centering
\small
\begin{tabular}{ll}
\toprule
\textbf{POS tag} & \textbf{Mean score} \\
\midrule
[CLS] & 1.00 (fixed) \\
Noun & 0.75 \\
Verb & 0.72 \\
Adjective & 0.82 \\
Adverb & 0.68 \\
Determiner & 0.25 \\
Preposition & 0.22 \\
Punctuation & 0.15 \\
\bottomrule
\end{tabular}
\end{table}

\textbf{Observation.} Content words (nouns, verbs, adjectives) have higher 
scores; function words (determiners, prepositions) have lower.

\subsubsection{Correlation with TF-IDF}

We compute the Pearson correlation between learned scores and TF-IDF:

\[
\rho_\text{Pearson} = \text{corr}(s, \text{TF-IDF}) = 0.65
\tag{8.1}
\]

\textbf{Interpretation.} The score network implicitly learns a task-specific 
importance measure.

\subsection{Per-Layer Analysis}

\textbf{Table 8.2: Effective kept ratio (DistilBERT).}

\begin{table}[H]
\centering
\small
\begin{tabular}{lll}
\toprule
\textbf{Layer} & \textbf{Kept ratio} & \textbf{$\Delta$ from target} \\
\midrule
0 & 85\% & 5\% \\
1 & 82\% & 2\% \\
2 & 80\% & 0\% \\
3 & 78\% & -2\% \\
4 & 77\% & -3\% \\
5 & 76\% & -4\% \\
\bottomrule
\end{tabular}
\end{table}

\textbf{Observation.} Later layers prune more aggressively.

\subsection{Comparison with Heuristic Pruning}

\textbf{Table 8.3: Learned vs heuristic scoring.}

\begin{table}[H]
\centering
\small
\begin{tabular}{llll}
\toprule
\textbf{Scoring} & \textbf{DistilBERT} & \textbf{RoBERTa} & \textbf{GPT-2} \\
\midrule
Attention magnitude & 85.5 & 73.0 & 74.0 \\
Token magnitude & 82.0 & 70.0 & 71.0 \\
\textbf{Learned MLP} & \textbf{89.5} & \textbf{77.0} & \textbf{78.0} \\
\bottomrule
\end{tabular}
\end{table}

\textbf{Observation.} Learned MLP improves by 3.5\%.

\subsection{Causality Verification}

We verify causality via finite differences:

\begin{lstlisting}
x = torch.randn(1, L, d, requires_grad=True)
out = model(x)
out[0, i].sum().backward()
assert x.grad[0, i+1:].abs().max() < 1e-6
\end{lstlisting}

\textbf{Result.}   Verified on all three models

\subsection{Effect of Score Network Depth}

\textbf{Table 8.4: Accuracy vs MLP depth.}

\begin{table}[H]
\centering
\small
\begin{tabular}{lll}
\toprule
\textbf{Depth} & \textbf{DistilBERT} & \textbf{Params/layer} \\
\midrule
Linear (1 layer) & 87.5 & ~768 \\
MLP-64 (2 layers) & \textbf{89.5} & ~50K \\
MLP-64-64 (3 layers) & 89.5 & ~55K \\
\bottomrule
\end{tabular}
\end{table}

\textbf{Observation.} One hidden layer suffices.

\subsection{Failure Cases}

\subsubsection{Cases Where PTP Fails}

\begin{table}[H]
\centering
\small
\begin{tabular}{ll}
\toprule
\textbf{Input type} & \textbf{Accuracy drop} \\
\midrule
Very short sequences ($L < 16$) & 5 \\
Adversarial inputs & 15 \\
Rare-word inputs & 3 \\
\bottomrule
\end{tabular}
\end{table}

\subsubsection{Root Causes}

Score network relies on content; adversarial inputs confuse it

\subsection{Numerical Stability}

\subsubsection{Score Clamping}

$10^4$ clamp prevents overflow. Score range is bounded.

\subsubsection{Gather Safety}

\texttt{topk} returns valid indices; gather is safe.

\subsubsection{Empty Mask}

If all tokens are padding, we clamp to first token.

\subsection{Summary of Insights}

\begin{enumerate}
\item \textbf{Learned scores correlate with content-word importance.}
\item \textbf{[CLS] preservation is essential.}
\item \textbf{Late layers prune more.}
\item \textbf{Learned MLP beats heuristics.}
\item \textbf{Causality preserved via position comparison.}
\end{enumerate}